% Capítulo 1: Prolegómenos
% Versión revisada tras la revisión del estado del arte

A continuación, se describe la motivación del presente proyecto y su alcance. También se incluye un glosario de términos y la organización del resto de la presente documentación.

\section{Motivación}

La comunicación es una de las actividades que definen a la especie humana. Ya que nos permite construir relaciones, expresar necesidades, compartir ideas y, en última instancia, participar en sociedad. Sin embargo, para un porcentaje significativo de la población, esta capacidad natural representa un desafío diario.

Las personas con Trastorno del Espectro Autista (TEA), entre otras condiciones, encuentran a menudo dificultades importantes para comunicarse verbalmente. Según la Organización Mundial de la Salud, aproximadamente 1 de cada 127 personas presenta autismo a nivel mundial~\cite{who2025autism}, y en España se estima que existen más de 450.000 personas con esta condición~\cite{autismoespana2024}. Estas condiciones hacen que su comunicación se vea afectada, y que su día a día esté profundamente condicionado.

¿Qué abordan estos sistemas tecnológicos versus generaciones humanas? Mediante los llamados Sistemas Aumentativos y Alternativos de Comunicación (SAAC). Estos sistemas permiten a las personas con dificultades del habla expresarse por vías alternativas, y entre los más utilizados están los basados en pictogramas: símbolos gráficos que representan conceptos, acciones u objetos de forma visual.

El problema es que las aplicaciones disponibles hoy en el mercado, aunque funcionales, tienen sus limitaciones. Herramientas como \textit{LetMeTalk}~\cite{letmetalk2024} o \textit{Proloquo2Go}~\cite{proloquo2go2024} ofrecen tableros de pictogramas con síntesis de voz básica, pero sus interfaces pueden resultar confusas y poco intuitivas. Aunque el problema va más allá de lo visual. Además, estas herramientas en muchos casos se dedican a concatenar palabras. Si un usuario selecciona los pictogramas ``yo'', ``querer'' y ``agua'', el sistema reproduce exactamente eso: ``yo querer agua''. Una construcción artificial, que obliga al interlocutor a hacer el trabajo de interpretación, y puede generar frustración en el emisor por no poder comunicarse como le gustaría.

Otro de los problemas de los sistemas actuales es el siguiente: mientras el habla natural alcanza entre 125 y 185 palabras por minuto, los sistemas AAC actuales raramente superan las 8-15~\cite{elsahar2019aac}. Esta diferencia de un orden de magnitud no es solo una cuestión de eficiencia; condiciona el tipo de comunicación posible, relegando a los usuarios a intercambios transaccionales cuando lo que necesitan es poder mantener conversaciones reales.

Más allá de lo técnico, este Trabajo de Fin de Grado parte de una premisa fundamental: la comunicación es fundamental en todo lo que hacemos. Es el medio a través del cual se construyen relaciones, se emprenden proyectos y se participa plenamente en la vida social. Poder comunicarse ---y hacerlo de forma efectiva--- no constituye un lujo, sino un elemento central del bienestar y la realización personal. Limitar esta capacidad en una persona supone restringir su desarrollo de formas que trascienden lo meramente práctico.

A estas limitaciones ya mencionadas se suma otra carencia significativa en las soluciones actuales: no aprovechan los avances recientes en inteligencia artificial. Ni los modelos de lenguaje que podrían generar frases naturales a partir de pictogramas, ni los motores de síntesis de voz neuronal que hoy pueden producir locuciones prácticamente indistinguibles de una voz humana.

Sin embargo, el momento actual resulta especialmente interesante para abordar estas carencias. Los modelos de lenguaje han alcanzado un nivel de madurez que permite generar texto natural de alta calidad, y la síntesis de voz ha experimentado avances comparables. Estas tecnologías, combinadas adecuadamente, ofrecen una oportunidad real de transformar la comunicación aumentativa.

En este contexto surge el presente proyecto, cuyo objetivo es desarrollar una aplicación que permita construir frases mediante pictogramas y convertirlas en habla natural. Para ello, se implementarán y compararán distintos enfoques de generación de lenguaje ---desde gramáticas formales hasta modelos de lenguaje de gran escala--- y se integrará un motor de síntesis de voz neuronal. El prototipo resultante se validará con usuarios reales para garantizar que el impacto trascienda el ámbito teórico.

\section{Alcance y objetivos}

El objetivo principal de este proyecto es desarrollar una aplicación multiplataforma que permita a personas con dificultades del habla construir frases seleccionando pictogramas y convertirlas en voz de forma natural. Para ello, el tiene los siguientes objetivos específicos:

\begin{enumerate}
    \item Análisis del estado del arte en sistemas SAAC y técnicas de generación de lenguaje natural aplicadas a comunicación aumentativa.
    \item Selección y preparación de una base de datos de pictogramas y frases frecuentes, orientada al dominio de comunicación cotidiana de personas con TEA.
    \item Implementación y evaluación comparativa de tres enfoques de IA para la generación de oraciones: gramáticas formales, representaciones mediante \textit{embeddings} y modelos de lenguaje de gran escala (LLM).
    \item Integración de un motor de síntesis de voz (TTS) que reproduzca las frases generadas de forma clara y natural.
    \item Diseño e implementación de una interfaz accesible e intuitiva, pensada específicamente para las necesidades de los usuarios objetivo.
    \item Validación del prototipo con usuarios reales.
    \item Documentación técnica y manual de uso.
\end{enumerate}

El sistema está pensado principalmente para personas con TEA que presentan dificultades en la comunicación verbal, sin restricción de edad. La validación con usuarios reales se realizará en colaboración con profesionales o entidades del ámbito (TODO: Incluir aquí las diversas entidades conforme vayamos contactando con ellas).

\section{Glosario de términos}

\begin{description}
    \item[ARASAAC] Centro Aragonés para la Comunicación Aumentativa y Alternativa. Ofrece pictogramas y recursos de forma gratuita bajo licencia Creative Commons. Es la referencia más utilizada en el mundo hispanohablante.
    
    \item[\textit{Embedding}] Representación de palabras o conceptos como vectores numéricos en un espacio donde elementos similares quedan cerca entre sí. Permite a los algoritmos trabajar con el significado de las palabras.
    
    \item[GLN / NLG] Generación de Lenguaje Natural (\textit{Natural Language Generation}). Subcampo del PLN enfocado en producir texto a partir de datos o representaciones estructuradas.
    
    \item[Gramática formal] Conjunto de reglas que definen cómo construir oraciones válidas en un lenguaje. En este proyecto, se usa para generar frases correctas a partir de pictogramas.
    
    \item[LLM (\textit{Large Language Model})] Modelo de lenguaje de gran escala, basado en redes neuronales, capaz de generar y comprender texto natural. Ejemplos conocidos son GPT o LLaMA.
    
    \item[Pictograma] Símbolo gráfico que representa un concepto, acción u objeto. Es el elemento básico de muchos sistemas de comunicación aumentativa.
    
    \item[PLN] Procesamiento del Lenguaje Natural. Área de la inteligencia artificial que trabaja con la comprensión y generación de lenguaje humano por parte de máquinas.
    
    \item[Realización superficial] (\textit{Surface realization}) Etapa final de la generación de lenguaje natural donde una representación abstracta se convierte en texto gramaticalmente correcto.
    
    \item[SAAC] Sistema Aumentativo y Alternativo de Comunicación. Herramientas y estrategias que facilitan la comunicación a personas con dificultades en el habla.
    
    \item[TEA] Trastorno del Espectro Autista. Condición del desarrollo que afecta a la comunicación e interacción social, con un amplio rango de manifestaciones.
    
    \item[TTS (\textit{Text-to-Speech})] Tecnología que convierte texto en voz sintetizada.
\end{description}

\section{Organización del documento}

\textit{TODO: Esta sección se completará conforme avance el documento.}